% TOP_PROBLEM: {"id": 2147, "title": "Erdős Problem #503", "category_id": 2, "category": {"id": 2, "name": "combinatorics", "display_name": "Combinatorics", "description": "Counting problems, graph theory, discrete structures.", "slug": "combinatorics", "order_index": 2, "created_at": "2026-07-31T15:26:25.671Z"}, "status": "open"}
% FIRST_POSTED: null
% ENTRY_KIND: statement_only
% Imported from: unsolvedmath_additional_500.tex; UnsolvedMath ID 2147
% !TeX program = lualatex
% Compile twice: lualatex 2147.tex
% Self-contained: no external bibliography, images, or \input files.
% Numeric section labels and visible numbers are original UnsolvedMath IDs.
\documentclass[11pt,a4paper]{article}
\usepackage[margin=24mm,headheight=14pt]{geometry}
\usepackage{fontspec}
\setmainfont{Latin Modern Roman}
\setsansfont{Latin Modern Sans}
\setmonofont{Latin Modern Mono}
\usepackage{amsmath,amssymb,mathtools}
\usepackage{unicode-math}
\setmathfont{Latin Modern Math}
\usepackage{microtype}
\usepackage{enumitem,xurl,needspace}
\usepackage[dvipsnames]{xcolor}
\usepackage{hyperref}
\hypersetup{unicode=true,colorlinks=true,linkcolor=MidnightBlue,urlcolor=MidnightBlue,
 pdftitle={UnsolvedMath Problem 2147},pdfauthor={Source-annotated compilation},bookmarksnumbered=true}
\usepackage{bookmark,tocloft,titlesec,fancyhdr}
\setlength{\cftsecnumwidth}{4.2em}
\cftsetpnumwidth{2.5em}
\cftsetrmarg{4em}
\titleformat{\section}{\Large\bfseries\raggedright}{\thesection}{0.7em}{}
\pagestyle{fancy}\fancyhf{}
\fancyhead[L]{\small\sffamily Additional UnsolvedMath statements}
\fancyhead[R]{\small Original catalogue IDs}
\fancyfoot[C]{\thepage}
\setlength{\parindent}{0pt}
\setlength{\parskip}{5pt plus 1pt}
\setlength{\emergencystretch}{3em}
\setcounter{tocdepth}{1}\setcounter{secnumdepth}{1}
\setlist[itemize]{leftmargin=3em,itemsep=4pt,topsep=4pt,parsep=0pt}
\allowdisplaybreaks
\newcommand{\N}{\mathbb{N}}
\newcommand{\Z}{\mathbb{Z}}
\newcommand{\Q}{\mathbb{Q}}
\newcommand{\R}{\mathbb{R}}
\newcommand{\C}{\mathbb{C}}
\newcommand{\F}{\mathbb{F}}
\newcommand{\A}{\mathbb{A}}
\newcommand{\PP}{\mathbb{P}}
\newcommand{\E}{\mathbb{E}}
\newcommand{\eps}{\varepsilon}
\newcommand{\dd}{\,\mathrm d}
\newcommand{\sref}[2]{\hyperlink{source:#1:#2}{[S#2]}}
\newif\ifshowcatalogue
\showcataloguetrue
\newcommand{\cataloguescope}[1]{\ifshowcatalogue
 \paragraph{Statement recorded in the supplied source.}
 #1\fi}
\begin{document}
% UnsolvedMath ID 2147; source catalogue code EP-503

\setcounter{section}{2146}

\section{Largest Euclidean Sets with Every Triple Isosceles}\label{2147}

{\small\sffamily UnsolvedMath ID 2147 · Combinatorics}

\subsection{Definitions and mathematical statement}

\paragraph{Shared notation.}Unless a section says otherwise, $\N=\{1,2,\ldots\}$,
$\Z$ is the ring of integers, $\Q,\R,\C$ are the rational, real, and complex
fields, $[N]=\{1,\ldots,\lfloor N\rfloor\}$, and $\log$ is the natural
logarithm. For real functions of a parameter tending to infinity,
$f\ll g$ and $f=O(g)$ mean $|f|\leq Cg$ eventually for some fixed $C$;
$f\gg g$ reverses this comparison; $f=o(g)$ means $f/g\to0$;
$f\sim g$ means $f/g\to1$; and $f\asymp g$ means both order bounds.
A subscript on $O$ or $\ll$ specifies allowed constant dependence.
The expression $n^{o(1)}$ in an upper bound means growth smaller than
$n^\epsilon$ for each fixed $\epsilon>0$. Local definitions govern any
exception, finite-field convention, or use of the same symbol for another
object. An asymptotic claim is not automatically an effective algorithm.



The displayed equal-distance condition is the precise convention, including any degenerate collinear triple satisfying it. The problem does not add noncollinearity. The ambient dimension is fixed when maximizing the finite cardinality.

Distances are Euclidean. Distinct distances mean distinct numerical values; when a distance is said to occur a certain number of times, unordered pairs of distinct points are counted. Congruence allows any Euclidean isometry, including reflection. \sref{2147}{1} \sref{2147}{2}

\paragraph{Statement recorded in the supplied source.}
What is the size of the largest $A\subseteq \mathbb{R}^d$ such that every three points from $A$ determine an isosceles triangle? That is, for any three points $x,y,z$ from $A$, at least two of the distances $\lvert x-y\rvert,\lvert y-z\rvert,\lvert x-z\rvert$ are equal. \sref{2147}{1}

\paragraph{Residual task reported by the archive.}

The embedded review (2026-08-17) records: Determine the extremum or close the quadratic-size gap in general dimension. \sref{2147}{1}

\subsection{Short English statement}

How many points can be placed in a fixed Euclidean dimension when every triple has two equal pairwise distances?

\subsection{Sources}

{\small

\begin{itemize}

\item[\hypertarget{source:2147:1}{S1}] ULAM.AI, \emph{UnsolvedMath}, supplied \texttt{problems.json}, record 2147 (EP-503), statement and background. The embedded literature-review date is 2026-08-17; that is the archive’s review date, not a fresh certification. Dataset landing page: \url{https://huggingface.co/datasets/ulamai/UnsolvedMath}.

\item[\hypertarget{source:2147:2}{S2}] T. F. Bloom, \emph{Erdős Problems}, Problem 503, \url{https://www.erdosproblems.com/503}. Reference imported from the supplied record; the live problem page was not independently checked for this compilation.

\item[\hypertarget{source:2147:3}{S3}] Z. Kovács, A note on Erdős's mysterious remark, arXiv:2412.05190 (2024). \url{https://arxiv.org/abs/2412.05190}. Bibliographic information imported from the supplied record.

\item[\hypertarget{source:2147:4}{S4}] Blokhuis, A., Few-distance sets. (1984), iv+70. Bibliographic information imported from the supplied record.

\item[\hypertarget{source:2147:5}{S5}] Croft, H. T., $9$-point and $7$-point configurations in $3$-space. Proc. London Math. Soc. (3) (1962), 400-424. Bibliographic information imported from the supplied record.

\item[\hypertarget{source:2147:6}{S6}] Erd\H{o}s, Paul and Kelly, L. M., Elementary Problems and Solutions: Solutions: E735. Amer. Math. Monthly (1947), 227-229. Bibliographic information imported from the supplied record.

\end{itemize}

}

\label{end:2147}
\end{document}
